\begin{lemma}
\label{lemma-invariance-number-branches-smooth}
Let $A \to B$ be a local homomorphism of local rings which
is the localization of a smooth ring map.
\begin{enumerate}
\item The number of geometric branches of $A$ is equal to the number
of geometric branches of $B$.
\item If $A \to B$ induces a purely inseparable extension of residue fields,
then the number of branches of $A$ is the number of branches of $B$.
\end{enumerate}
\end{lemma}

\begin{proof}
We will use that smooth ring maps are flat
(Algebra, Lemma \ref{algebra-lemma-smooth-syntomic}),
that localizations are flat (Algebra, Lemma
\ref{algebra-lemma-flat-localization}),
that compositions of flat ring maps are flat
(Algebra, Lemma \ref{algebra-lemma-composition-flat}),
that base change of a flat ring map is flat
(Algebra, Lemma \ref{algebra-lemma-flat-base-change}),
that flat local homomorphisms are faithfully flat
(Algebra, Lemma \ref{algebra-lemma-local-flat-ff}),
that (strict) henselization is flat
(Lemma \ref{lemma-dumb-properties-henselization}), and
Going down for flat ring maps
(Algebra, Lemma \ref{algebra-lemma-flat-going-down}).

\medskip\noindent
Proof of (2). Let $A^h$, $B^h$ be the henselizations of $A$, $B$. Then $B^h$
is the henselization of $A^h \otimes_A B$ at the unique
maximal ideal lying over $\mathfrak m_B$, see
Algebra, Lemma \ref{algebra-lemma-henselian-functorial-improve}.
Thus we may and do assume $A$ is henselian.
Since $A \to B \to B^h$ is flat, every minimal prime of $B^h$
lies over a minimal prime of $A$ and since $A \to B^h$ is faithfully flat,
every minimal prime of $A$ does lie under a minimal prime of $B^h$;
in both cases use going down for flat ring maps.
Therefore it suffices to show that given a minimal prime
$\mathfrak p \subset A$, there is at most one minimal prime
of $B^h$ lying over $\mathfrak p$.
After replacing $A$ by $A/\mathfrak p$ and $B$ by $B/\mathfrak p B$
we may assume that $A$ is a domain; the $A$ is still henselian by
Algebra, Lemma \ref{algebra-lemma-quotient-henselization}.
By Lemma \ref{lemma-unibranch} we see that the integral closure
$A'$ of $A$ in its field of fractions is a local domain.
Of course $A'$ is a normal domain. By
Algebra, Lemma \ref{algebra-lemma-normal-goes-up}
we see that $A' \otimes_A B^h$ is a normal ring
(the lemma just gives it for $A' \otimes_A B$, to go up to
$A' \otimes_A B^h$ use that $B^h$ is a colimit of \'etale
$B$-algebras and use Algebra, Lemma \ref{algebra-lemma-colimit-normal-ring}).
By Algebra, Lemma \ref{algebra-lemma-local-tensor-with-integral}
we see that $A' \otimes_A B^h$ is local (this is where we
use the assumption on the residue fields of $A$ and $B$).
Hence $A' \otimes_A B^h$ is a local normal ring, hence a local domain.
Since $B^h \subset A' \otimes_A B^h$ by flatness of $A \to B^h$
we conclude that $B^h$ is a domain as desired.

\medskip\noindent
Proof of (1). Let $A^{sh}$, $B^{sh}$ be strict henselizations
of $A$, $B$. Then $B^{sh}$ is a strict henselization of
$A^h \otimes_A B$ at a maximal ideal lying over $\mathfrak m_B$
and $\mathfrak m_{A^h}$, see
Algebra, Lemma \ref{algebra-lemma-strictly-henselian-functorial-improve}.
Thus we may and do assume $A$ is strictly henselian.
Since $A \to B \to B^{sh}$ is flat, every minimal prime of $B^{sh}$
lies over a minimal prime of $A$ and since $A \to B^{sh}$ is faithfully flat,
every minimal prime of $A$ does lie under a minimal prime of $B^{sh}$;
in both cases use going down for flat ring maps.
Therefore it suffices to show that given a minimal prime
$\mathfrak p \subset A$, there is at most one minimal prime
of $B^{sh}$ lying over $\mathfrak p$.
After replacing $A$ by $A/\mathfrak p$ and $B$ by $B/\mathfrak p B$
we may assume that $A$ is a domain; then $A$ is still strictly henselian by
Algebra, Lemma \ref{algebra-lemma-quotient-strict-henselization}.
By Lemma \ref{lemma-geometrically-unibranch} we see that the integral closure
$A'$ of $A$ in its field of fractions is a local domain whose residue
field is a purely inseparable extension of the residue field of $A$.
Of course $A'$ is a normal domain. By
Algebra, Lemma \ref{algebra-lemma-normal-goes-up}
we see that $A' \otimes_A B^{sh}$ is a normal ring
(the lemma just gives it for $A' \otimes_A B$, to go up to
$A' \otimes_A B^{sh}$ use that $B^{sh}$ is a colimit of \'etale
$B$-algebras and use Algebra, Lemma \ref{algebra-lemma-colimit-normal-ring}).
By Algebra, Lemma \ref{algebra-lemma-local-tensor-with-integral}
we see that $A' \otimes_A B^{sh}$ is local (since $A \subset A'$
induces a purely inseparable residue field extension).
Hence $A' \otimes_A B^{sh}$ is a local normal ring, hence a local domain.
Since $B^{sh} \subset A' \otimes_A B^{sh}$ by flatness of $A \to B^{sh}$
we conclude that $B^{sh}$ is a domain as desired.
\end{proof}